\section{Contracting forced differences before Euler optimization}
\label{sec:cocycle-quotient}

The exact Euler optimizer of Section~\ref{sec:cocycle-euler} still routes
flow through vertices whose relative potentials cannot change anywhere on
the minimum-span face. The new implementation identifies these forced
differences, solves a smaller problem and reconstructs a full certificate
on the original indexed model. The independent arithmetic checker and all
source-bound surface verifiers remain unchanged. If no vertices can be
contracted, an adaptive branch retains the original full-network solver.

\subsection{Which tight constraints force equality}

Start with the verified feasible potential $p^0$. For a difference constraint
$p_b-p_a\leq d_{ab}$, its slack is
$r_{ab}=d_{ab}+p^0_a-p^0_b\geq0$. Put $q_v=p_v-p^0_v$, so the constraint becomes
$q_b-q_a\leq r_{ab}$. On a directed zero-slack path, the shifted potential
cannot increase. If two vertices are mutually reachable using such paths,
their shifted potentials must be equal for \emph{every} feasible solution.
Thus each strongly connected component of the zero-slack graph is one
quotient variable. A one-way tight edge is insufficient: $q_b\leq q_a$
still allows strict inequality. The implementation uses iterative Kosaraju
traversals, avoiding recursion depth proportional to the graph size.

For a component map $C$, write $p_v=p^0_v+q_{C(v)}$. An objective term becomes
\[
 w_e\left|c_e+p^0_b-p^0_a+q_{C(b)}-q_{C(a)}\right|.
\]
Internal terms are constants and need no flow optimization. Between two
components, terms with the same oriented endpoints and shifted offset are
combined by adding their weights; reversing an orientation also negates the
offset. Among constraints with the same ordered component pair, only the
smallest shifted bound is needed. Internal constraints are redundant because
their shifted left side is zero and their bounds are nonnegative.

These operations preserve the complete feasible set after the forced
substitution and preserve the objective up to its explicitly retained
constant terms. They do not choose one representative surface topology or
restrict the face to sampled extrema. The zero quotient potential is already
feasible, so the existing exact flow solver applies directly to the reduced
network. The final potential is expanded and normalized only by an additive
constant. Its objective is recomputed on the original edge list.

\subsection{Lifting the dual without trusting the quotient}

The returned certificate still has one flow per original objective edge and
per original constraint. An internal objective edge receives the sign-saturating
flow $w_e\operatorname{sign}(c_e+p^0_b-p^0_a)$. A zero-weight edge receives zero.
For an aggregated intercomponent term, split its signed flow $F\in[-W_g,W_g]$
across the original intervals $[-w_e,w_e]$: begin at all lower endpoints and
distribute $F+W_g$ units in bulk, at most $2w_e$ to each member. Restore each
member's original orientation. This uses one pass over the members and never
iterates over flow units. Because all members have the same shifted affine
expression, the split preserves complementary equality. Each quotient
constraint flow is assigned to one original constraint attaining its minimum
bound; the discarded parallel constraints receive zero.

The lifted flow now balances on each component as a whole, but may not balance
at its individual vertices. Two directed trees inside the zero-slack component
repair this. A tree directed toward the root comes from the reverse traversal;
a forward traversal supplies a tree directed away from it. In reverse order
on the inward tree, send each positive vertex imbalance toward the root,
allowing intervening negative imbalances to absorb flow. Every nonroot vertex
then has nonpositive imbalance. In reverse order on the outward tree, send
the remaining subtree demand from parent to child. Component balance ensures
that the root also ends balanced.

All added constraint flows are nonnegative and use original zero-slack arcs.
Their constraints stay tight after the quotient potential is expanded, so
they preserve complementary equality. Internal constant objective terms,
aggregated terms and retained constraint flows therefore give a feasible dual
on the original model with the same value as its expanded primal potential.
The unchanged checker independently verifies every original flow bound,
difference inequality, vertex balance equation and objective equality. It
receives neither a component partition nor a compressed certificate whose
interpretation would depend on the producer.

\subsection{Operation bounds and an adaptive no-contraction branch}

Let $n,e,m$ be the original counts of vertices, objective terms and constraints;
let $k,e',m'$ be their quotient counts, and let $W'$ be the remaining objective
weight sum. Construction, expansion and tree routing require
$O(n+e+m)$ indexed graph operations, apart from dictionary operations used to
combine rows. The reduced shortest-path phase costs
\[
 O\!\left(W'(k+e'+m')\log(k+1)\right)
\]
arithmetic and comparison operations. For the geometric model,
$n+e+m=O(N)$ and $W'\leq W\leq12N$, retaining the previous worst-case polynomial
bound. There is no universal asymptotic improvement when the quotient is as
large as the original graph. If $k=n$, the implementation bypasses aggregation
and lifting and calls the original solver; the additional cost is the linear
zero-slack graph analysis. Python hash-collision accounting and binary integer
costs remain subject to the conservative polynomial qualifications of the
preceding section.

Lifting is independent of the numerical magnitude of its flows. The reduced
solver sends at most $W'$ units along simple paths with $O(k)$ arcs, so the sum
of its final nonnegative constraint flows is $O(kW')$. Together with the bounded
objective flows, the initial positive imbalance to repair is at most
$O(W+kW')$. Each tree edge carries at most this amount per traversal. Thus
all lifted flows have polynomial magnitude in the geometric input size and
$O(\log(N+1))$ bits. Their sum need not be linear in $N$. The
solver's auxiliary capacity $W+1$ is not a constraint on the original dual:
lifted constraint flows need only be nonnegative.
Heights and shifted costs remain binary encoded; no level, sheet or flow unit
is expanded. Independent arithmetic replay remains linear in the original model size
in this model; source and geometry checks are retained.

A concrete family separates the costs. Take $n$ vertices in a directed cycle
of zero-slack constraints and one unit-weight objective edge from vertex zero
to every other vertex, each with shifted offset one. The old initialization
creates $n-1$ separate unit supplies, so the full solver performs $n-1$
augmentations. Each shortest-path call scans the source's $n-1$ arc entries,
giving a quadratic scan lower bound for that implementation. The new quotient
has one vertex, no variable objective terms and no augmentations; two tree
passes lift its dual in linear graph work. The retained operation-count
experiment measures this family directly, rather than inferring an
asymptotic improvement from wall-clock ratios alone.

\subsection{Preserved optima and a newly completed large source}

The comparison loads the full flow implementation from revision
\code{18e76fe87}. Its complete source hash is retained. Before comparing the
new implementation, it exactly reproduces every source record, positive proof,
arithmetic dual and Regina result in the preceding audit. All 83 previously
completed sources preserve their Euler optimum, minimum piece count and
independently checked connected-surface summary under contraction. Their
selected potentials need not be the same: 58 source-certificate hashes change.
Every resulting source surface and arithmetic dual is checked again.

The new run completes 84 of 85 sources under the same shared two-million-guard
allowance. Gordian remains capped during the unchanged span solver. Circle 33,
previously capped in Euler optimization, now finishes the entire audited
pipeline at 1,197,679 guards; its Euler stage takes 90,106 guards. Its face
contracts from 130 variables to 64. This is a newly completed exact-face
experiment, not a newly recognized knot: the default native search already
has an earlier witness for this source.

All completed sources have a nontrivial contraction. Selected reductions are
14 to 4 variables for the trefoil, 22 to 2 for the original-disc fixture,
30 to 5 for the tied-disc fixture, 46 to 8 for random case 19, 50 to 13 for
random case 158, and 46 to 5 for Kinoshita--Terasaka. The corresponding Euler
stage guard counts fall from 23,470 to 5,909; 58,288 to 9,579; 108,324 to 13,618;
292,201 to 21,604; 428,401 to 25,311; and 322,516 to 21,253. Guard counts measure
the deterministic allowance, not elapsed time or bit operations.

The audit retains 74 distinct full arithmetic models and lifted duals and 14
distinct positive source proofs, covering 19 positive source labels. Regina
reconfirms all 75 small-source controls; their observed boundary counts also
agree with the previous optima. Neither that agreement nor equality of Euler
characteristic is used to assert equality of all possible surface topologies.
The certified negative optimum for the three-crossing unknot remains $-1$.
The restriction to the minimum-span face and the incomplete nature of general
surface discovery are unchanged.

The forced-cycle family gives the following measured arithmetic guard and arc
scan counts. Every quotient has one vertex and no variable objective terms.
\begin{center}
\small
\begin{tabular}{rrrrr}
\toprule
$n$ & Full guards & New guards & Full scans & New augmentations \\
\midrule
8 & 503 & 150 & 231 & 0 \\
16 & 1,851 & 302 & 975 & 0 \\
32 & 7,043 & 606 & 3,999 & 0 \\
64 & 27,411 & 1,214 & 16,191 & 0 \\
128 & 108,083 & 2,430 & 65,151 & 0 \\
\bottomrule
\end{tabular}

\end{center}
The new flow stage uses $19n-2$ measured guards at these sizes; the preceding argument
establishes linear indexed graph work without depending on an empirical fit.
Additional tests cover 240 random forced-group problems against the full
solver, reversed and parallel objective terms, tight one-way constraints,
malformed lifted duals, cancellation and the adaptive singleton bypass. A
1,500-vertex cycle with 20,000-bit vertex offsets checks iterative traversal,
large cancellation in the dual, binary JSON replay and a linear guard bound.
The earlier 240 exhaustive-grid arithmetic controls remain active.

\subsection{Complete candidate-pipeline measurements}

All 18 focused tests pass in 1.905 seconds, and all 1,223 maintained tests
pass in 223.627 seconds. The source comparison and cycle-family audit takes
35.539 seconds. The correctness audit and full suite partly overlap; these
elapsed values are not paired performance measurements. Both the audit and
benchmark pin 505 source files. The arithmetic dual checker, public Euler
model construction, native recognizer and existing independent positive
verifiers remain byte-identical to the preceding revision.

The isolated timing experiment compares the old and new flow functions
within the same complete candidate pipeline: source diagram, finite exterior,
primitive cocycle, span optimization, Euler optimization, independent
source-bound surface replay and certificate hash. Both arms include the same
function-selection wrapper. The baseline module is loaded before timing.
There is no local work cap and no early native witness short-circuit; these
are candidate-pipeline measurements rather than full knot recognition.

Five rounds randomize old, old A/A, new and new A/A calls for eight sources,
following one warm-up round each. All 160 measured calls complete, with 32
warm-ups; the run takes 45.746 seconds. No tests, audits, builds or commits
overlap it. Per-arm Euler values, piece counts, operation counts and hashes
are stable across rounds. The primary-arm medians and median paired ratios
are as follows; the ratios need not equal quotients of the displayed medians.
\begin{center}
\small
\begin{tabular}{lrrrrr}
\toprule
Input & Full (ms) & Quotient (ms) & Full/new & Full A/A & New A/A \\
\midrule
Trefoil & 35.807 & 31.535 & 1.128 & 1.001 & 1.006 \\
Original disc & 68.724 & 57.351 & 1.196 & 0.996 & 1.001 \\
Tied disc & 107.483 & 85.417 & 1.253 & 0.999 & 0.984 \\
Random 19 & 204.390 & 135.757 & 1.506 & 1.014 & 0.993 \\
Random 158 & 245.875 & 151.331 & 1.625 & 1.013 & 0.996 \\
Kinoshita--Terasaka & 207.558 & 134.843 & 1.539 & 1.017 & 1.011 \\
Circle 17 & 352.748 & 243.128 & 1.442 & 1.005 & 0.978 \\
Circle 33 & 1072.704 & 643.779 & 1.670 & 0.989 & 1.017 \\
\bottomrule
\end{tabular}

\end{center}
All eight selected pipelines improve, by paired factors from 1.128 to 1.670.
The A/A medians range from 0.978 to 1.017. The smaller complete-pipeline gains,
compared with the reduction in Euler guard work, include unchanged source,
cohomology, span and verification costs. This evidence does not establish a
uniform speedup on arbitrary indexed models: when no forced equality exists,
the new preprocessing is additional linear work before the full solver.

Contraction is now used inside the explicit Euler optimizer, with the
singleton fallback described above. The default native recognition schedule
remains unchanged, and there is no claim of a default-portfolio speedup.
The update improves exact optimization within the existing restricted family
and gives an asymptotic improvement on the forced-cycle family; it does not
supply a general subexponential or quasipolynomial knot-recognition bound.
Reproduce the audit and timings with \code{normal\_orbit\_research.quotient},
in \code{audit} or \code{benchmark --rounds 5} mode. The checked incoming
unknot inventory and \code{openai/math} head \code{fd4aeeb2e} remain unchanged.


